\chapter{Analyse et Spécification des besoins}

\section*{Introduction}
Dans ce chapitre, nous identifierons les acteurs du projet et définirons les besoins fonctionnels et non fonctionnels. Ensuite, nous détaillerons toutes les fonctionnalités du projet à l'aide d'un diagramme de cas d'utilisation global et de ses raffinements. Enfin, nous présenterons le planning sous la forme d'un diagramme de Gantt et décrirons l'approche Scrum adoptée pour la gestion du projet.

\section{Spécification des besoins}
Dans cette section, nous présenterons les différents acteurs de l'application, ainsi que les besoins fonctionnels et non fonctionnels identifiés pour le projet.

\subsection{Identification des acteurs}
L'acteur, qu'il soit interne ou externe, physique ou logiciel, est une entité capable d'interagir directement avec le système applicatif. Dans ce cadre, nous avons identifié quatre acteurs, comme présenté dans le tableau \ref{tab:acteurIdent} ci-dessous :

\begin{longtable}[c]{|p{0.3\linewidth}|p{0.65\linewidth}|}
  \hline
  \textbf{Acteurs}            & \textbf{Description}                                                                                                                                                                                                                                                                                           \\
  \hline
  \endfirsthead

  \hline
  \textbf{Acteurs}            & \textbf{Description}                                                                                                                                                                                                                                                                                           \\
  \hline
  \endhead

  \hline
  \endfoot

  \hline
  \endlastfoot

  Administrateur              & Dispose d'un accès aux fonctionnalités d'administration de la plateforme. Il peut gérer les utilisateurs, attribuer les rôles, activer ou désactiver les comptes, créer et supprimer des espaces de travail, inviter des membres et supprimer des projets.                                                     \\
  \hline

  Chef de projet              & Responsable de la gestion des projets et de la collaboration au sein des espaces de travail auxquels il appartient. Il peut consulter les espaces accessibles, créer et modifier des projets, inviter des membres, gérer les tâches, le backlog, les sprints, les réunions et suivre l'avancement du projet.   \\
  \hline

  Membre                      & Représente un utilisateur appartenant à un espace de travail ou à un projet. Il peut consulter les projets auxquels il est affecté, participer aux discussions, consulter et mettre à jour les tâches, commenter les éléments du backlog, suivre les sprints et accéder aux réunions selon ses droits d'accès. \\
  \hline

  Utilisateur non authentifié & Peut accéder aux interfaces d'inscription et de connexion afin de créer un compte ou de s'authentifier pour accéder aux fonctionnalités de l'application.                                                                                                                                                      \\
  \hline

  \caption{\centering Identification des Acteurs et leurs Fonctions dans le Système Applicatif}
  \label{tab:acteurIdent}
\end{longtable}

\underline{\textbf{Remarque}}: Les fonctionnalités accessibles dépendent du rôle de l'utilisateur ainsi que de son appartenance à un espace de travail ou à un projet.

\subsection{Besoins fonctionnels}

La mise en \oe{}uvre de la future application Planora doit répondre aux besoins fonctionnels suivants :
\begin{itemize}
  \item[$\bullet$] \textbf{Gestion d'inscription :}
        \begin{itemize}
          \item[$\circ$] Création d'un nouveau compte utilisateur.
          \item[$\circ$] Enregistrement des informations personnelles de l'utilisateur.
        \end{itemize}

  \item[$\bullet$] \textbf{Gestion de l'authentification :}
        \begin{itemize}
          \item[$\circ$] Connexion sécurisée à l'application.
          \item[$\circ$] Gestion de la session utilisateur.
        \end{itemize}

  \item[$\bullet$] \textbf{Gestion des utilisateurs et de leurs droits d'accès :}
        \begin{itemize}
          \item[$\circ$] Consultation de la liste des utilisateurs.
          \item[$\circ$] Attribution des rôles aux utilisateurs.
          \item[$\circ$] Activation et désactivation des comptes utilisateurs.
        \end{itemize}

  \item[$\bullet$] \textbf{Gestion des espaces de travail :}
        \begin{itemize}
          \item[$\circ$] Création et consultation des espaces de travail.
          \item[$\circ$] Invitation des membres à rejoindre un espace de travail.
          \item[$\circ$] Acceptation ou refus des invitations reçues.
          \item[$\circ$] Suppression des membres d'un espace de travail.
        \end{itemize}

  \item[$\bullet$] \textbf{Gestion des projets :}
        \begin{itemize}
          \item[$\circ$] Création, consultation, modification et suppression des projets.
          \item[$\circ$] Ajout et suppression des membres d'un projet.
          \item[$\circ$] Invitation des membres à rejoindre un projet.
        \end{itemize}

  \item[$\bullet$] \textbf{Gestion du backlog :}
        \begin{itemize}
          \item[$\circ$] Création, modification et suppression des éléments du backlog.
          \item[$\circ$] Affectation des éléments du backlog aux membres.
          \item[$\circ$] Définition de la priorité, de la complexité, des story points et de la date d'échéance.
          \item[$\circ$] Ajout de commentaires, de liens, de branches, de commits et de pièces jointes.
        \end{itemize}

  \item[$\bullet$] \textbf{Gestion des tâches et sous-tâches :}
        \begin{itemize}
          \item[$\circ$] Création, consultation, modification et suppression des tâches.
          \item[$\circ$] Affectation des tâches aux membres du projet.
          \item[$\circ$] Suivi de l'état d'avancement des tâches et des sous-tâches.
        \end{itemize}

  \item[$\bullet$] \textbf{Gestion des sprints :}
        \begin{itemize}
          \item[$\circ$] Création, consultation, modification et suppression des sprints.
          \item[$\circ$] Association des éléments du backlog à un sprint.
          \item[$\circ$] Suivi de l'avancement des sprints à travers un tableau de bord agile.
        \end{itemize}

  \item[$\bullet$] \textbf{Gestion de la collaboration :}
        \begin{itemize}
          \item[$\circ$] Échange de messages entre les membres d'un projet.
          \item[$\circ$] Consultation des conversations liées aux projets.
          \item[$\circ$] Épinglage des messages importants.
        \end{itemize}

  \item[$\bullet$] \textbf{Gestion des réunions :}
        \begin{itemize}
          \item[$\circ$] Création et consultation des réunions.
          \item[$\circ$] Planification des réunions dans un calendrier.
          \item[$\circ$] Définition des membres concernés par chaque réunion.
        \end{itemize}

  \item[$\bullet$] \textbf{Suivi et analyse de la santé de l'équipe :}
        \begin{itemize}
          \item[$\circ$] Analyse des messages échangés entre les membres.
          \item[$\circ$] Détection du stress, de la frustration et de l'ambiance générale.
          \item[$\circ$] Génération d'indicateurs et d'alertes liés à la santé de l'équipe.
        \end{itemize}

  \item[$\bullet$] \textbf{Consultation du tableau de bord :}
        \begin{itemize}
          \item[$\circ$] Visualisation de l'avancement des projets.
          \item[$\circ$] Consultation des statistiques globales et des indicateurs de suivi.
        \end{itemize}
\end{itemize}

\subsection{Besoins non fonctionnels}
Les besoins non fonctionnels sont tout aussi cruciaux que les besoins fonctionnels dans tout projet, affectant la qualité, la sécurité et l'expérience utilisateur globale. Ils sont essentiels pour garantir la performance et la satisfaction des utilisateurs finaux.

\begin{itemize}[label=\textbullet]
  \item \textbf{Ergonomie :} L'application doit disposer d'une interface ergonomique, intuitive et conforme aux bonnes pratiques UI/UX, afin d'offrir une expérience utilisateur optimale.

  \item \textbf{Fiabilité :} L'application doit garantir une disponibilité élevée, des temps de réponse réduits et une gestion efficace des erreurs.

  \item \textbf{Maintenabilité :} Le code doit être organisé de manière claire, modulaire et documentée, afin de faciliter sa compréhension, son évolution et sa maintenance par d'autres développeurs.

  \item \textbf{Sécurité :} L'application doit intégrer un dispositif de sécurité cohérent garantissant l'intégrité, la disponibilité et la confidentialité des données, ainsi que l'authentification des utilisateurs.
\end{itemize}

\section {Diagramme de cas d'utilisation global}
Le diagramme de cas d'utilisation global présenté dans la figure \ref{fig:Use Case} offre une vue d'ensemble des interactions des acteurs avec les diverses fonctionnalités de l'application cible, permettant ainsi une compréhension rapide des fonctionnalités et des besoins des utilisateurs.

\begin{figure}[H]
  \centering
  \includegraphics[width=1\linewidth]{UseCaseDiagram1.png}
  \caption{Diagramme de cas d'utilisation globale}
  \label{fig:Use Case}
\end{figure}

\section{Raffinement des cas d'utilisation}
\subsection{Gestion des membres et authentification}

 %La figure \ref{fig:GestionUtilisateurUseCase} ci-dessous illustre le diagramme de raffinement de la gestion des membres et de l'authentification.

 %\begin{figure}[H]
  % \centering
   % \includegraphics[width=1\columnwidth]{gererM.png}
   %\caption{Diagramme de cas d'utilisation « Gestion des utilisateurs et authentification »}
   % %\label{fig:GestionUtilisateurUseCase}
 %\end{figure}

Les tableaux ci-dessous présentent une description textuelle des cas d'utilisation liés à la gestion des utilisateurs.

\subsubsection{Utilisateur}

\begin{longtable}[c]{|p{0.3\linewidth}|p{0.65\linewidth}|}
  \hline
  \multicolumn{2}{|c|}{\textbf{S'authentifier}}                                                                                                  \\
  \hline
  \endfirsthead

  \hline
  \multicolumn{2}{|c|}{\textbf{S'authentifier}}                                                                                                  \\
  \hline
  \endhead

  \hline
  \endfoot

  \hline
  \endlastfoot

  \textbf{Objectif}      & Permettre à l'utilisateur d'accéder à l'application Planora selon son rôle.                                           \\
  \hline
  \textbf{Acteur(s)}     & Administrateur, ProjectManager, Membre                                                                                \\
  \hline
  \textbf{Précondition}  & L'utilisateur doit posséder un compte valide.                                                                         \\
  \hline
  \textbf{Scénario}      &
  1. L'utilisateur accède à la page de connexion.                                                                                                \\
                         & 2. L'utilisateur saisit son adresse email et son mot de passe.                                                        \\
                         & 3. Le système vérifie les informations d'identification.                                                              \\
                         & 4. Si les informations sont valides, le système génère une session et redirige l'utilisateur vers le tableau de bord. \\
                         & 5. Si les informations sont incorrectes, le système affiche un message d'erreur.                                      \\
  \hline
  \textbf{Postcondition} & L'utilisateur est authentifié et accède aux fonctionnalités autorisées selon son rôle.                                \\
  \hline
  \caption{\centering Description textuelle du cas d'utilisation « S'authentifier »}
\end{longtable}

\subsubsection{Administrateur}

\begin{longtable}[c]{|p{0.3\linewidth}|p{0.65\linewidth}|}
  \hline
  \multicolumn{2}{|c|}{\textbf{Gérer les utilisateurs}}                                                                                 \\
  \hline
  \endfirsthead

  \hline
  \multicolumn{2}{|c|}{\textbf{Gérer les utilisateurs}}                                                                                 \\
  \hline
  \endhead

  \hline
  \endfoot

  \hline
  \endlastfoot

  \textbf{Objectif}      & Administrer les comptes utilisateurs et gérer leurs rôles dans l'application.                                \\
  \hline
  \textbf{Acteur(s)}     & Administrateur                                                                                               \\
  \hline
  \textbf{Précondition}  & L'administrateur doit être authentifié.                                                                      \\
  \hline
  \textbf{Scénario}      &
  1. L'administrateur accède à la liste des utilisateurs.                                                                               \\
                         & 2. Le système affiche les comptes enregistrés dans l'application.                                            \\
                         & 3. L'administrateur sélectionne un utilisateur.                                                              \\
                         & 4. L'administrateur peut modifier ses informations, lui attribuer un rôle, activer ou désactiver son compte. \\
                         & 5. Le système enregistre les modifications effectuées.                                                       \\
  \hline
  \textbf{Postcondition} & Les informations ou les droits d'accès de l'utilisateur sont mis à jour.                                     \\
  \hline
  \caption{\centering Description textuelle du cas d'utilisation « Gérer les utilisateurs »}
\end{longtable}
\subsection{Gestion des espaces de travail}

La figure \ref{fig:GestionWorkspaceUseCase} ci-dessous illustre le diagramme de raffinement de la gestion des espaces de travail.

\begin{figure}[H]
  \centering
  \includegraphics[width=1\columnwidth]{gereE.png}
  \caption{Diagramme de cas d'utilisation « Gestion des espaces de travail »}
  \label{fig:GestionWorkspaceUseCase}
\end{figure}

\begin{longtable}[c]{|p{0.3\linewidth}|p{0.65\linewidth}|}
  \hline
  \multicolumn{2}{|c|}{\textbf{Gérer un espace de travail}}                                                             \\
  \hline
  \endfirsthead

  \hline
  \multicolumn{2}{|c|}{\textbf{Gérer un espace de travail}}                                                             \\
  \hline
  \endhead

  \hline
  \endfoot

  \hline
  \endlastfoot

  \textbf{Objectif}      & Créer, consulter et organiser les espaces de travail de l'application.                       \\
  \hline
  \textbf{Acteur(s)}     & Administrateur, ProjectManager, Membre                                                       \\
  \hline
  \textbf{Précondition}  & L'utilisateur doit être authentifié.                                                         \\
  \hline
  \textbf{Scénario}      &
  1. L'utilisateur accède à la section des espaces de travail.                                                          \\
                         & 2. Le système affiche les espaces de travail accessibles.                                    \\
                         & 3. L'administrateur peut créer un nouvel espace de travail.                                  \\
                         & 4. L'administrateur ou le ProjectManager peut inviter des membres selon les droits accordés. \\
                         & 5. Le membre invité peut accepter ou refuser l'invitation reçue.                             \\
                         & 6. Le système met à jour la liste des membres de l'espace de travail.                        \\
  \hline
  \textbf{Postcondition} & L'espace de travail est créé ou mis à jour avec les membres concernés.                       \\
  \hline
  \caption{\centering Description textuelle du cas d'utilisation « Gérer un espace de travail »}
\end{longtable}
\subsection{Gestion des projets}

La figure \ref{fig:GestionProjetUseCase} ci-dessous illustre le diagramme de raffinement de la gestion des projets.

\begin{figure}[H]
  \centering
   \includegraphics[width=1\columnwidth]{gereP.png}
  \caption{Diagramme de cas d'utilisation « Gestion des projets »}
  \label{fig:GestionProjetUseCase}
\end{figure}

\begin{longtable}[c]{|p{0.3\linewidth}|p{0.65\linewidth}|}
  \hline
  \multicolumn{2}{|c|}{\textbf{Gérer un projet}}                                                             \\
  \hline
  \endfirsthead

  \hline
  \multicolumn{2}{|c|}{\textbf{Gérer un projet}}                                                             \\
  \hline
  \endhead

  \hline
  \endfoot

  \hline
  \endlastfoot

  \textbf{Objectif}      & Permettre la création, la modification, la suppression et le suivi des projets.                   \\
  \hline
  \textbf{Acteur(s)}     & Administrateur, ProjectManager, Membre                                            \\
  \hline
  \textbf{Précondition}  & L'utilisateur doit être authentifié et appartenir à l'espace de travail concerné. \\
  \hline
  \textbf{Scénario}      &
  1. L'utilisateur accède à la liste des projets.                                                            \\
                         & 2. Le système affiche les projets accessibles.                                    \\
                         & 3. Le ProjectManager crée ou modifie un projet.                                   \\
                         & 4. Le ProjectManager invite des membres à rejoindre le projet.                    \\
                         & 5. Les membres acceptent ou refusent l'invitation.                                \\
                         & 6. Le système met à jour les informations du projet et sa liste des membres.      \\
  \hline
  \textbf{Postcondition} & Le projet est créé, modifié ou mis à jour avec les membres affectés.              \\
  \hline
  \caption{\centering Description textuelle du cas d'utilisation « Gérer un projet »}
\end{longtable}
\subsection{Gestion du backlog, des tâches et des sprints}

La figure \ref{fig:GestionBacklogSprintUseCase} ci-dessous illustre le diagramme de raffinement de la gestion du backlog, des tâches et des sprints.

\begin{figure}[H]
  \centering
  \includegraphics[width=1\columnwidth]{gereB.png}
  \caption{Diagramme de cas d'utilisation « Gestion du backlog, des tâches et des sprints »}
  \label{fig:GestionBacklogSprintUseCase}
\end{figure}

\begin{longtable}[c]{|p{0.3\linewidth}|p{0.65\linewidth}|}
  \hline
  \multicolumn{2}{|c|}{\textbf{Gérer le backlog et les sprints}}                                                                                              \\
  \hline
  \endfirsthead

  \hline
  \multicolumn{2}{|c|}{\textbf{Gérer le backlog et les sprints}}                                                                                              \\
  \hline
  \endhead

  \hline
  \endfoot

  \hline
  \endlastfoot

  \textbf{Objectif}      & Organiser les éléments du backlog, les tâches et les sprints d'un projet.                                                          \\
  \hline
  \textbf{Acteur(s)}     & ProjectManager, Membre                                                                                                             \\
  \hline
  \textbf{Précondition}  & L'utilisateur doit être authentifié et avoir accès au projet concerné.                                                             \\
  \hline
  \textbf{Scénario}      &
  1. L'utilisateur accède au backlog du projet.                                                                                                               \\
                         & 2. Le système affiche les éléments du backlog existants.                                                                           \\
                         & 3. L'utilisateur peut créer, modifier ou supprimer un élément du backlog.                                                          \\
                         & 4. L'utilisateur peut affecter un élément à un membre, définir sa priorité, sa complexité, ses story points et sa date d'échéance. \\
                         & 5. L'utilisateur crée un sprint et y associe des éléments du backlog.                                                              \\
                         & 6. Le système met à jour l'état d'avancement des tâches et des sprints.                                                            \\
  \hline
  \textbf{Postcondition} & Le backlog, les tâches ou les sprints sont mis à jour selon les actions effectuées.                                                \\
  \hline
  \caption{\centering Description textuelle du cas d'utilisation « Gérer le backlog et les sprints »}
\end{longtable}
\subsection{Gestion de la collaboration et des réunions}

La figure \ref{fig:GestionCollaborationUseCase} ci-dessous illustre le diagramme de raffinement de la gestion de la collaboration et des réunions.

\begin{figure}[H]
  \centering
   \includegraphics[width=1\columnwidth]{gereR.png}
  \caption{Diagramme de cas d'utilisation « Gestion de la collaboration et des réunions »}
  \label{fig:GestionCollaborationUseCase}
\end{figure}

\begin{longtable}[c]{|p{0.3\linewidth}|p{0.65\linewidth}|}
  \hline
  \multicolumn{2}{|c|}{\textbf{Collaborer au sein d'un projet}}                                                               \\
  \hline
  \endfirsthead

  \hline
  \multicolumn{2}{|c|}{\textbf{Collaborer au sein d'un projet}}                                                               \\
  \hline
  \endhead

  \hline
  \endfoot

  \hline
  \endlastfoot

  \textbf{Objectif}      & Faciliter la communication entre les membres d'un projet et organiser les réunions.                \\
  \hline
  \textbf{Acteur(s)}     & ProjectManager, Membre                                                                             \\
  \hline
  \textbf{Précondition}  & L'utilisateur doit être authentifié et appartenir au projet concerné.                              \\
  \hline
  \textbf{Scénario}      &
  1. L'utilisateur accède à l'espace de discussion du projet.                                                                 \\
                         & 2. Le système affiche les conversations associées au projet.                                       \\
                         & 3. L'utilisateur envoie ou consulte les messages échangés avec les membres.                        \\
                         & 4. L'utilisateur peut consulter les réunions planifiées dans le calendrier.                        \\
                         & 5. Le ProjectManager peut créer une réunion et définir les membres concernés.                      \\
                         & 6. Le système enregistre la réunion et l'affiche dans le calendrier.                               \\
  \hline
  \textbf{Postcondition} & Les messages ou les réunions du projet sont enregistrés et consultables par les membres concernés. \\
  \hline
  \caption{\centering Description textuelle du cas d'utilisation « Collaborer au sein d'un projet »}
\end{longtable}
\subsection{Suivi de la santé de l'équipe}

La figure \ref{fig:TeamHealthUseCase} ci-dessous illustre le diagramme de raffinement du suivi de la santé de l'équipe.

\begin{figure}[H]
  \centering
   \includegraphics[width=1\columnwidth]{ConsAi.png}
  \caption{Diagramme de cas d'utilisation « Suivi de la santé de l'équipe »}
  \label{fig:TeamHealthUseCase}
\end{figure}

\begin{longtable}[c]{|p{0.3\linewidth}|p{0.65\linewidth}|}
  \hline
  \multicolumn{2}{|c|}{\textbf{Analyser la santé de l'équipe}}                                                           \\
  \hline
  \endfirsthead

  \hline
  \multicolumn{2}{|c|}{\textbf{Analyser la santé de l'équipe}}                                                           \\
  \hline
  \endhead

  \hline
  \endfoot

  \hline
  \endlastfoot

  \textbf{Objectif}      & Analyser les messages échangés afin d'évaluer l'ambiance générale de l'équipe.                \\
  \hline
  \textbf{Acteur(s)}     & ProjectManager, Membre                                                                        \\
  \hline
  \textbf{Précondition}  & L'utilisateur doit être authentifié et avoir accès au projet concerné.                        \\
  \hline
  \textbf{Scénario}      &
  1. L'utilisateur accède à la page de suivi de la santé de l'équipe.                                                    \\
                         & 2. Le système récupère les messages récents échangés dans le projet.                          \\
                         & 3. Le module d'intelligence artificielle analyse les messages.                                \\
                         & 4. Le système classe les messages selon leur sentiment : positif, neutre, stressé ou frustré. \\
                         & 5. Le système calcule un score global et affiche les indicateurs de santé de l'équipe.        \\
                         & 6. En cas de stress ou de frustration élevés, le système affiche des alertes.                 \\
  \hline
  \textbf{Postcondition} & Les indicateurs et alertes liés à la santé de l'équipe sont affichés à l'utilisateur.         \\
  \hline
  \caption{\centering Description textuelle du cas d'utilisation « Analyser la santé de l'équipe »}
\end{longtable}
\subsection{Communication entre les membres d'un projet}

La figure \ref{fig:TeamHealthUseCase} ci-dessous illustre le diagramme de raffinement du communication entre les memebres d'un projet.

\begin{figure}[H]
  \centering
   \includegraphics[width=1\columnwidth]{ComAi.png}
  \caption{Diagramme de cas d'utilisation « Communication entre les membres d'un projet »}
  \label{fig:TeamHealthUseCase}
\end{figure}

\begin{longtable}[c]{|p{0.3\linewidth}|p{0.65\linewidth}|}
  \hline
  \multicolumn{2}{|c|}{\textbf{Communication entre les membres d'un projet}}                                                           \\
  \hline
  \endfirsthead

  \hline
  \multicolumn{2}{|c|}{\textbf{Communication entre les membres d'un projet}}                                                           \\
  \hline
  \endhead

  \hline
  \endfoot

  \hline
  \endlastfoot

  \textbf{Objectif}      & Analyser les messages échangés afin d'évaluer l'ambiance générale de l'équipe.                \\
  \hline
  \textbf{Acteur(s)}     & Admin, ProjectManager, Membre                                                                        \\
  \hline
  \textbf{Précondition}  & L'utilisateur doit être authentifié et avoir accès au projet concerné.                        \\
  \hline
  \textbf{Scénario}      &
  1. L'utilisateur accède à la page de suivi de la santé de l'équipe.                                                    \\
                         & 2. Le système récupère les messages récents échangés dans le projet.                          \\
                         & 3. Le module d'intelligence artificielle analyse les messages.                                \\
                         & 4. Le système classe les messages selon leur sentiment : positif, neutre, stressé ou frustré. \\
                         & 5. Le système calcule un score global et affiche les indicateurs de santé de l'équipe.        \\
                         & 6. En cas de stress ou de frustration élevés, le système affiche des alertes.                 \\
  \hline
  \textbf{Postcondition} & Les indicateurs et alertes liés à la santé de l'équipe sont affichés à l'utilisateur.         \\
  \hline
  \caption{\centering Description textuelle du cas d'utilisation « Analyser la santé de l'équipe »}
\end{longtable}


\section*{Conclusion}


Dans le deuxième chapitre, nous avons nous avons effectué une analyse et une spécification des besoins en identifiant les acteurs clés du projet.Ensuite nous avons défini les exigences fonctionnelles et non fonctionnelles. Enfin, nous avons présenté une illustration du diagramme d'utilisation global ainsi que ses raffinements.
